% Generated draft: requires mathematical and source review.
\documentclass[11pt]{article}
\usepackage[margin=25mm]{geometry}
\usepackage{fontspec}
\setmainfont{Latin Modern Roman}
\usepackage{amsmath,amssymb,mathtools,unicode-math}
\setmathfont{Latin Modern Math}
\usepackage{enumitem,needspace,xurl,hyperref}
\setlength{\parindent}{0pt}
\setlength{\parskip}{5pt}
\setlength{\emergencystretch}{3em}
\newcommand{\N}{\mathbb{N}}
\newcommand{\Z}{\mathbb{Z}}
\newcommand{\Q}{\mathbb{Q}}
\newcommand{\R}{\mathbb{R}}
\newcommand{\C}{\mathbb{C}}
\newcommand{\F}{\mathbb{F}}
\newcommand{\A}{\mathbb{A}}
\newcommand{\PP}{\mathbb{P}}
\newcommand{\E}{\mathbb{E}}
\newcommand{\eps}{\varepsilon}
\newcommand{\dd}{\,\mathrm{d}}
\begin{document}
\section*{UnsolvedMath ID: 30005656}
\textbf{Source problem code:} catalogue entry
\subsection{Definitions and mathematical statement}
Let \(p\) be a prime, including \(p=2\), and let \(G=C_p\) be the cyclic group of order \(p\). Let \(R\) be a commutative ring and write \(\underline R\) for its constant \(G\)-Mackey functor: its value at every orbit is \(R\), restrictions and conjugations are identities, and the transfer associated with \(H\leq K\leq G\) is multiplication by \([K:H]\). Write \(H\underline R\) for the associated genuine \(G\)-equivariant Eilenberg--Mac Lane ring spectrum, denoted \(HR\) in the source.

Under the usual interpretation of the source notation, let
\[
\mathcal T_R=\operatorname{Perf}_{C_p}(H\underline R)
\]
be the tensor triangulated category of compact genuine \(C_p\)-equivariant \(H\underline R\)-modules, with tensor product over \(H\underline R\). The supplied excerpt does not explicitly define its convention for \(\operatorname{Perf}_{C_p}\); this interpretation requires confirmation against the report.

For an essentially small tensor triangulated category \(\mathcal T\), its Balmer spectrum \(\operatorname{Spc}(\mathcal T)\) consists of its prime thick tensor ideals. A thick tensor ideal \(\mathfrak p\) is a full triangulated subcategory closed under direct summands and under tensoring with every object of \(\mathcal T\). It is prime if it is proper and
\[
X\otimes Y\in\mathfrak p\quad\Longrightarrow\quad X\in\mathfrak p\ \text{or}\ Y\in\mathfrak p.
\]
The Balmer topology has the supports
\[
\operatorname{supp}(X)=\{\mathfrak p\in\operatorname{Spc}(\mathcal T):X\notin\mathfrak p\}
\]
as a closed subbasis.

\textbf{Problem (source Question 8).} For every prime \(p\) and every commutative ring \(R\), describe the topological space
\[
\operatorname{Spc}\bigl(\operatorname{Perf}_{C_p}(H\underline R)\bigr).
\]
The source supplies the case in which \(p\) is invertible in \(R\) as
\[
p\in R^\times\quad\Longrightarrow\quad
\operatorname{Spc}(\mathcal T_R)\cong\operatorname{Spec}(R).
\]
Thus its remaining question concerns \(p\notin R^\times\). Here \(\operatorname{Spec}(R)\) is the prime ideal spectrum with its Zariski topology.

\textbf{Supplied structural context.} For a subgroup \(H\leq G\), write \(\Phi^H\) for geometric fixed points and \(W_G(H)=N_G(H)/H\) for the Weyl group. The source attributes to Barthel--Castellana--Heard--Naumann-Pol, under its reference label [1], a description for finite groups of the form
\[
\operatorname{Spc}\bigl(\operatorname{Perf}_{G}(H\underline R)\bigr)
\simeq
\coprod_{(H)\leq G}
\operatorname{Spc}\bigl(\operatorname{Perf}(\Phi^H H\underline R)\bigr)/W_G(H)
\qquad\text{as sets},
\]
where \((H)\) ranges over conjugacy classes of subgroups, and each quotient is an orbit set. For \(G=C_p\), with \(e\) the trivial subgroup, the supplied excerpt reduces this to
\[
\operatorname{Spc}(\mathcal T_R)
\simeq
\operatorname{Spc}\bigl(\operatorname{Perf}(\Phi^{C_p}H\underline R)\bigr)
\amalg
\operatorname{Spc}\bigl(\operatorname{Perf}(\Phi^eH\underline R)\bigr)/C_p.
\]
It further states that \(H\underline R\) is a global spectrum and that the indicated \(C_p\)-action is trivial, giving a set bijection
\[
\operatorname{Spc}(\mathcal T_R)
\simeq
\operatorname{Spc}\bigl(\operatorname{Perf}(\Phi^{C_p}H\underline R)\bigr)
\amalg
\operatorname{Spc}\bigl(\operatorname{Perf}(\Phi^eH\underline R)\bigr).
\]
These are supplied set descriptions, not assertions that the Balmer topology is the disjoint union topology. The requested description must determine the topology, including the specialization relations between the two displayed parts. The precise hypotheses and construction underlying the source's global-spectrum and trivial-action assertions require review.
\subsection{Short English statement}
For each prime \(p\) and commutative ring \(R\), describe the prime tensor ideals of perfect genuine \(C_p\)-equivariant modules over the Eilenberg--Mac Lane spectrum of the constant Mackey functor \(\underline R\), together with their Balmer topology. The source supplies the answer when \(p\) is invertible and a geometric fixed point description of the underlying set. Its remaining target is the topology and how the parts fit together when \(p\) is not invertible. The supplied open status and structural assertions require human mathematical and source review.
\subsection{Sources}
\begin{itemize}
\item Supplied local UnsolvedMath registry record, permanent ID \(30005656\), title \emph{Balmer Spectra for Constant Mackey Functors}, exact source problem code \texttt{catalogue entry}, in the Oberwolfach Reports Open Problems collection. The record assigns proposed year 2023 and supplies no named proposer.
\item Supplied report citation: \emph{Tensor-Triangular Geometry and Interactions (2024)}, Oberwolfach Reports, pp.\ 2303--2358, DOI \texttt{10.4171/owr/2023/40}. \url{https://doi.org/10.4171/owr/2023/40}. The excerpt contains Question 8, Remark 9, and Theorem 10. The citation is supplied authoritative provenance; it has not been independently consulted for this draft.
\item Supplied literature source: \emph{Quillen stratification in equivariant homotopy theory}. \url{https://doi.org/10.1007/s00222-024-01301-0}.
\item Supplied literature source: \emph{Balmer and Gallauer on permutation modules}. \url{https://math.ucla.edu/~balmer/Pubfile/PmMfAm.pdf}.
\item The excerpt credits Balmer and Gallauer with describing the spectrum for finite \(G\) and fields \(k\) of characteristic \(p>0\), and with describing the \(C_2\), integral-coefficient case in previous work on Artin--Tate motives, using \(C_2\) as the Galois group. That previous work is cited there as [2], without its bibliographic entry in the supplied record. The geometric fixed point set description is attributed there to Barthel--Castellana--Heard--Naumann-Pol as Theorem 10, citing [1]; the record does not explicitly match [1] and [2] to complete bibliography entries.
\item Literature assessment dated exactly \texttt{2026-08-22}: the supplied curation labels the problem open and states that several field and integral coefficient cases are known, while the topology and gluing for general commutative rings are not fully determined. This is supplied dated curation derived from the OWR record, not an independently verified literature assessment.
\item The registry labels the statement \texttt{corrected\_verified} and reports that the official report was checked directly. These are supplied verification metadata, not verification performed for this draft.
\end{itemize}
\end{document}
